\section{Compatibilité avec un compilateur C (partie optionnelle)}
\label{sec:compilateur}

Jusqu'ici, les programmes étaient écrits à la main en hexadécimal. Le sujet demande les
étapes nécessaires pour que le processeur exécute le code produit par un compilateur
RISC-V, en considérant le compilateur \og RISC-V (32-bits) gcc (trunk)\fg{} du site Compiler
Explorer (\texttt{godbolt.org}). Nous les avons mises en œuvre à partir du processeur
synthétisable de la partie~\ref{sec:fpga} ; les fichiers sont dans \texttt{exo/p7}.

\subsection{Étapes nécessaires}

Le code produit par le compilateur suppose un processeur conforme à la spécification et un
environnement qui respecte la convention d'appel \texttt{ilp32} (pile, registres
d'arguments et de retour). Ce n'était pas le cas de notre système :
\begin{enumerate}
\item \textbf{Adressage par octets} : \texttt{pc} avance de 4 à chaque instruction et les
adresses de \texttt{lw} et \texttt{sw} sont des adresses d'octets, alors que notre modèle
comptait en mots de 32 bits (section~\ref{sec:adapter}).
\item \textbf{Jeu d'instructions RV32I complet} : le compilateur emploie librement
\texttt{lui}, \texttt{auipc}, \texttt{jal}, \texttt{jalr}, les six branchements, les
comparaisons, les décalages par une constante, etc. Il faut aussi interdire les
instructions compressées de 16 bits, en compilant avec \texttt{-march=rv32im} au lieu de
la valeur par défaut \texttt{rv32gc}.
\item \textbf{Code de démarrage et pile} : avant \texttt{main}, il faut initialiser le
pointeur de pile \texttt{sp} sur une mémoire de données, appeler \texttt{main}, puis arrêter
le processeur. Le résultat de \texttt{main} est rendu dans \texttt{a0} (x10).
\item \textbf{Édition de liens et chargement} : placer le code à l'adresse 0, résoudre les
appels de fonctions, puis produire le fichier hexadécimal lu par \texttt{\$readmemh}.
\item \textbf{Répartition des mémoires} : le code reste dans la ROM du processeur, la pile
est dans \texttt{quadram}.
\end{enumerate}

\subsection{Compilation sur Compiler Explorer}

Le programme de test (code~\ref{lst:programme-c}) utilise une boucle, un tableau local, deux
fonctions, une multiplication et un reste ; il n'a ni variable globale ni appel de
bibliothèque. Il est compilé avec \texttt{-O0 -march=rv32im -mabi=ilp32} : sans
optimisation, le compilateur ne précalcule pas le résultat et toutes les variables passent
par la pile, ce qui teste les accès à \texttt{quadram}. Le code assembleur obtenu
(figure~\ref{fig:p7-godbolt}) compte 104 instructions.

\noindent\begin{minipage}{\linewidth}
\begin{lstlisting}[style=code, caption={Programme C de test (\texttt{programme.c}, extrait)}, label={lst:programme-c}]
int factorielle(int n) { ... }          /* produit 2 x 3 x ... x n */
int pgcd(int a, int b) { ... }          /* algorithme d'Euclide, avec % */

int main(void)
{
    int carres[10];
    int somme = 0;
    for (int i = 0; i < 10; i++)
        carres[i] = i * i;
    for (int i = 0; i < 10; i++)
        somme += carres[i];                              /* 285 */
    return somme + factorielle(6) + pgcd(1071, 462);    /* 285 + 720 + 21 */
}
\end{lstlisting}
\end{minipage}\par

\capture{p7-godbolt.png}{Capture de godbolt.org : \texttt{programme.c}, compilateur
RISC-V (32-bits) gcc (trunk), options \texttt{-O0 -march=rv32im -mabi=ilp32}}
{Compilation de \texttt{programme.c} sur Compiler Explorer}{fig:p7-godbolt}

\subsection{Assemblage et édition de liens}

Compiler Explorer fournit le code assembleur, mais pas un exécutable adapté à notre
processeur. Nous avons écrit \texttt{asm\_rv32.py}, qui joue le rôle de l'assembleur et de
l'éditeur de liens. En deux passes, il calcule l'adresse des étiquettes puis code chaque
instruction ; il comprend les noms de registres de l'ABI (\texttt{sp}, \texttt{a0}...) et
les pseudo-instructions de \texttt{gcc} (\texttt{li}, \texttt{mv}, \texttt{j}, \texttt{ble},
\texttt{call}...). Le fichier \texttt{crt0.s} (code~\ref{lst:crt0}) est placé à l'adresse 0,
avant le code compilé.

\noindent\begin{minipage}{\linewidth}
\begin{lstlisting}[style=code, caption={Code de démarrage \texttt{crt0.s}}, label={lst:crt0}]
_start:
    li     sp, 0x4000      # stack at the top of the quadram: 4 x 1024 words = 16 KB
    call   main            # main(); the result is returned in a0
    ebreak                 # stop: the processor is halted
\end{lstlisting}
\end{minipage}\par

Pour vérifier l'assembleur, nous avons aussi demandé à Compiler Explorer le code objet du
programme. Les 104 mots que nous produisons sont identiques à ceux du compilateur, sauf
les deux \texttt{jalr} des appels de fonctions : le compilateur y laisse un déplacement nul,
que l'éditeur de liens doit compléter, alors que \texttt{asm\_rv32.py} calcule le
déplacement réel. Le fichier final, \texttt{programme.txt}, compte 108 mots : 4 pour le
démarrage, 104 pour le programme.

\subsection{Modifications du processeur}

Le processeur \texttt{riscvproc\_rv32.v} reprend la machine à états de la
partie~\ref{sec:fpga}, avec les changements suivants :
\begin{itemize}
\item \texttt{pc} est une adresse d'octet qui avance de 4, et la ROM est lue avec
\texttt{pc[11:2]} ; l'adresse envoyée à \texttt{quadram} est l'adresse d'octet divisée par 4
(bits 13 à 2), ce qui donne toujours un mot de 12 bits ;
\item les instructions ajoutées sont celles du tableau~\ref{tab:rv32i-ajout} ;
\item \texttt{ecall} arrête le processeur comme \texttt{ebreak}.
\end{itemize}
Le décalage arithmétique est calculé à part, pour la raison vue à la
section~\ref{sec:rv32im-impl} : dans un opérateur \texttt{? :} dont une branche est non signée,
\lstinline[style=vlinline]|$signed(a) >>> n| deviendrait un décalage logique.

\begin{table}[H]
\centering
\begin{tabular}{ll}
\toprule
Format & Instructions ajoutées \\
\midrule
U & \texttt{lui}, \texttt{auipc} \\
J & \texttt{jal} \\
I & \texttt{jalr}, \texttt{slti}, \texttt{sltiu}, \texttt{xori}, \texttt{slli}, \texttt{srli}, \texttt{srai} \\
B & \texttt{beq}, \texttt{bne}, \texttt{blt}, \texttt{bge}, \texttt{bltu}, \texttt{bgeu} \\
R & \texttt{slt}, \texttt{sltu}, \texttt{xor}, \texttt{sra}, \texttt{or}, \texttt{and} \\
\bottomrule
\end{tabular}
\caption{Instructions ajoutées au processeur pour le compilateur}
\label{tab:rv32i-ajout}
\end{table}

Deux limites restent. \texttt{quadram} n'écrit que des mots entiers : les accès à un octet
ou à un demi-mot (\texttt{lb}, \texttt{sb}, \texttt{lh}, \texttt{sh}...), utilisés pour les
types \texttt{char} et \texttt{short}, sont refusés et signalés par la sortie
\texttt{illegal}. Par ailleurs, les variables globales initialisées devraient être recopiées
de la ROM vers \texttt{quadram} par le code de démarrage, ce que notre \texttt{crt0.s} ne
fait pas.

\subsection{Validation}

En simulation (figure~\ref{fig:p7-programme}), le programme compilé s'arrête sur
\texttt{ebreak} après 9\,910~ns, soit environ 990 périodes, avec \texttt{a0} = 1026
($285 + 720 + 21$). Le pointeur de pile est revenu à \texttt{0x4000} et \texttt{ra} vaut 12,
l'adresse de retour de l'appel à \texttt{main}. Aucune instruction n'a été refusée.

Le programme C n'utilise qu'une partie des instructions ajoutées. Le programme assembleur
\texttt{test\_rv32i.s} vérifie les autres en huit étapes : \texttt{lui}, \texttt{auipc},
comparaisons signées et non signées, opérations logiques, décalages logiques et
arithmétiques, branchements signés et non signés, appel et retour de fonction, et deux mots
écrits à 4 octets d'écart. Il se termine avec \texttt{a0} = 0, ce qui signifie que toutes
les vérifications sont passées (figure~\ref{fig:p7-test}).

\capture{p7-programme.png}{Transcript ModelSim, dossier \texttt{exo/p7},
\texttt{systeme\_rv32\_tb} avec \texttt{programme.txt} : registres à la fin}
{Exécution du programme C compilé : \texttt{a0} = 1026}{fig:p7-programme}

\capture{p7-test.png}{Transcript ModelSim, dossier \texttt{exo/p7},
\texttt{systeme\_rv32\_tb} avec \texttt{test\_rv32i.txt}}
{Programme de vérification des instructions ajoutées : \texttt{a0} = 0}{fig:p7-test}

Le processeur reste synthétisable. Compilé pour la même carte, avec le même affichage, il
utilise 6\,631 éléments logiques (64~\%), 20 blocs M9K (43~\%) et 8 multiplieurs ; sa
fréquence maximale est de 5,53~MHz, et il fonctionne lui aussi à 3,125~MHz. Sur la carte, les
afficheurs montrent 1026 ; KEY0 et KEY1 affichent \texttt{sp} = 16384 (\texttt{0x4000}) et
\texttt{ra} = 12, et RESET remet les registres à 0 (figure~\ref{fig:p7-carte}). Le code
produit par \texttt{gcc} s'exécute donc sur notre processeur.

\begin{figure}[H]
\centering
\photo[0.48\linewidth]{p7_1026.jpg}{sans touche : \texttt{a0} = 1026}\hfill
\photo[0.48\linewidth]{p7_16384.jpg}{KEY0 : \texttt{sp} = 16384}\\[3mm]
\photo[0.48\linewidth]{p7_12.jpg}{KEY1 : \texttt{ra} = 12}\hfill
\photo[0.48\linewidth]{p7_0.jpg}{RESET enfoncé : 0}
\caption{Le programme C exécuté sur la carte}
\label{fig:p7-carte}
\end{figure}
